% GENERATED FILE. Edit canonical lessons, then run npm run generate:books.
% canonical-source-hash: fnv1a64:488600a273c137c4
% canonical-lessons: GU-C321-talvar, GU-C321-dhal, GU-C321-hathiyar, GU-C321-top, GU-C321-vidhvu

\chapter{Swords and Shields}
\label{ch:gu-swords-and-shields}
\hlchaptermodality{321}

\begin{chapteropening}
\textbf{By the end of this chapter:} \emph{I can name the weapons of an old war and say how they were used.}

Complete the last lesson of chapter 321: I can name the weapons of an old war and say how they were used.
\end{chapteropening}

\section[\texorpdfstring{\gu{તલવાર} (talvār)}{તલવાર (talvār)}]{\gu{તલવાર} (talvār) — a sword}

\label{lesson:GU-C321-talvar}



\begin{quote}
Before the new one: say the Gujarati for an enemy, then the Gujarati for an invasion.
\end{quote}

\subsection*{You'll want to know: \gu{તલવાર}}
\textbf{\gu{તલવાર}} — \emph{talvār} — \textquotedblleft{}a sword\textquotedblright{}.

\begin{quote}
\gu{સૈનિકના} \gu{હાથમાં} \gu{તલવાર} \gu{હતી}.
\end{quote}

\begin{quote}
\emph{sainiknā hāthmā̃ talvār hatī.}
\end{quote}

\begin{quote}
The soldier had a sword in his hand.
\end{quote}

A feminine noun.

\begin{cousinweb}
\textbf{\gu{તલવાર}} (\emph{talvār}) goes back to Sanskrit \emph{taravāri}, \textquotedblleft{}a sword\textquotedblright{}, and CDIAL lists Gujarati \emph{tarvār}, \textquotedblleft{}a one-edged sword\textquotedblright{} (CDIAL 5706); the \emph{r} has become \emph{l} in the spoken word.
\end{cousinweb}

\begin{grammarlens}[title={What you've built}]
One more word for telling what happened.
\end{grammarlens}

\subsection*{Your turn}
\begin{itemize}
\raggedright
  \item \emph{Say it:} \emph{talvār}
  \item \emph{Say it:} \emph{talvār}, once more
  \item \emph{Say it:} \emph{sainiknā hāthmā̃ talvār hatī.} — The soldier had a sword in his hand.
  \item \emph{Recall:} say the Gujarati for an enemy, then the Gujarati for an invasion, then say \emph{talvār} again
\end{itemize}

\begin{tcolorbox}[breakable,colback=teal!4,colframe=teal!35!black,title={Before you move on}]
What does \gu{તલવાર} mean? (\textquotedblleft{}A sword\textquotedblright{}.) Say it once more.
\end{tcolorbox}
\section[\texorpdfstring{\gu{ઢાલ} (ḍhāl)}{ઢાલ (ḍhāl)}]{\gu{ઢાલ} (ḍhāl) — a shield}

\label{lesson:GU-C321-dhal}



\begin{quote}
Before the new one: say the Gujarati for an invasion, then the Gujarati for a sword.
\end{quote}

\subsection*{You'll want to know: \gu{ઢાલ}}
\textbf{\gu{ઢાલ}} — \emph{ḍhāl} — \textquotedblleft{}a shield\textquotedblright{}.

\begin{quote}
\gu{તેણે} \gu{ઢાલથી} \gu{પોતાને} \gu{બચાવ્યો}.
\end{quote}

\begin{quote}
\emph{teṇe ḍhālthī potāne bacāvyo.}
\end{quote}

\begin{quote}
He protected himself with his shield.
\end{quote}

A feminine noun.

\begin{cousinweb}
\textbf{\gu{ઢાલ}} (\emph{ḍhāl}) goes back to \emph{ḍhāla}, \textquotedblleft{}a shield\textquotedblright{}, and CDIAL lists the Gujarati noun (CDIAL 5583). It is feminine: \textbf{\gu{ઢાલથી}}, with a shield.
\end{cousinweb}

\begin{grammarlens}[title={What you've built}]
One more word for telling what happened.
\end{grammarlens}

\subsection*{Your turn}
\begin{itemize}
\raggedright
  \item \emph{Say it:} \emph{ḍhāl}
  \item \emph{Say it:} \emph{ḍhāl}, once more
  \item \emph{Say it:} \emph{teṇe ḍhālthī potāne bacāvyo.} — He protected himself with his shield.
  \item \emph{Recall:} say the Gujarati for an invasion, then the Gujarati for a sword, then say \emph{ḍhāl} again
\end{itemize}

\begin{tcolorbox}[breakable,colback=teal!4,colframe=teal!35!black,title={Before you move on}]
What does \gu{ઢાલ} mean? (\textquotedblleft{}A shield\textquotedblright{}.) Say it once more.
\end{tcolorbox}
\section[\texorpdfstring{\gu{હથિયાર} (hathiyār)}{હથિયાર (hathiyār)}]{\gu{હથિયાર} (hathiyār) — a weapon}

\label{lesson:GU-C321-hathiyar}



\begin{quote}
Before the new one: say the Gujarati for a sword, then the Gujarati for a shield.
\end{quote}

\subsection*{You'll want to know: \gu{હથિયાર}}
\textbf{\gu{હથિયાર}} — \emph{hathiyār} — \textquotedblleft{}a weapon\textquotedblright{}.

\begin{quote}
\gu{દુશ્મને} \gu{હથિયાર} \gu{નીચે} \gu{મૂક્યાં}.
\end{quote}

\begin{quote}
\emph{duśmane hathiyār nīce mūkyā̃.}
\end{quote}

\begin{quote}
The enemy laid down their weapons.
\end{quote}

A neuter noun.

\begin{cousinweb}
\textbf{\gu{હથિયાર}} (\emph{hathiyār}) comes from \emph{*hastakāra}, a form CDIAL reconstructs on \emph{hásta}, the hand, where it lists Gujarati \emph{hathiyār}, \textquotedblleft{}a tool\textquotedblright{} (CDIAL 14027): what the hand works with. In war it is a weapon.
\end{cousinweb}

\begin{grammarlens}[title={What you've built}]
One more word for telling what happened.
\end{grammarlens}

\subsection*{Your turn}
\begin{itemize}
\raggedright
  \item \emph{Say it:} \emph{hathiyār}
  \item \emph{Say it:} \emph{hathiyār}, once more
  \item \emph{Say it:} \emph{duśmane hathiyār nīce mūkyā̃.} — The enemy laid down their weapons.
  \item \emph{Recall:} say the Gujarati for a sword, then the Gujarati for a shield, then say \emph{hathiyār} again
\end{itemize}

\begin{tcolorbox}[breakable,colback=teal!4,colframe=teal!35!black,title={Before you move on}]
What does \gu{હથિયાર} mean? (\textquotedblleft{}A weapon\textquotedblright{}.) Say it once more.
\end{tcolorbox}
\section[\texorpdfstring{\gu{તોપ} (top)}{તોપ (top)}]{\gu{તોપ} (top) — a cannon}

\label{lesson:GU-C321-top}



\begin{quote}
Before the new one: say the Gujarati for a shield, then the Gujarati for a weapon.
\end{quote}

\subsection*{You'll want to know: \gu{તોપ}}
\textbf{\gu{તોપ}} — \emph{top} — \textquotedblleft{}a cannon\textquotedblright{}.

\begin{quote}
\gu{કિલ્લા} \gu{પર} \gu{બે} \gu{જૂની} \gu{તોપ} \gu{છે}.
\end{quote}

\begin{quote}
\emph{killā par be jūnī top che.}
\end{quote}

\begin{quote}
There are two old cannons on the fort.
\end{quote}

A feminine noun.

\begin{cousinweb}
\textbf{\gu{તોપ}} (\emph{top}) is usually traced to Turkish \emph{top}, \textquotedblleft{}a ball\textquotedblright{}, and then \textquotedblleft{}a cannon\textquotedblright{}, a word that came with the armies of the Mughal age. It is feminine: \textbf{\gu{જૂની} \gu{તોપ}}.
\end{cousinweb}

\begin{grammarlens}[title={What you've built}]
One more word for telling what happened.
\end{grammarlens}

\subsection*{Your turn}
\begin{itemize}
\raggedright
  \item \emph{Say it:} \emph{top}
  \item \emph{Say it:} \emph{top}, once more
  \item \emph{Say it:} \emph{killā par be jūnī top che.} — There are two old cannons on the fort.
  \item \emph{Recall:} say the Gujarati for a shield, then the Gujarati for a weapon, then say \emph{top} again
\end{itemize}

\begin{tcolorbox}[breakable,colback=teal!4,colframe=teal!35!black,title={Before you move on}]
What does \gu{તોપ} mean? (\textquotedblleft{}A cannon\textquotedblright{}.) Say it once more.
\end{tcolorbox}
\section[\texorpdfstring{\gu{વીંધવું} (vī̃dhvũ)}{વીંધવું (vī̃dhvũ)}]{\gu{વીંધવું} (vī̃dhvũ) — to pierce}

\label{lesson:GU-C321-vidhvu}



\begin{quote}
Before the new one: say the Gujarati for a weapon, then the Gujarati for a cannon.
\end{quote}

\subsection*{You'll want to know: \gu{વીંધવું}}
\textbf{\gu{વીંધવું}} — \emph{vī̃dhvũ} — \textquotedblleft{}to pierce\textquotedblright{}.

\begin{quote}
\gu{સૈનિકે} \gu{તલવારથી} \gu{ઢાલ} \gu{વીંધી}.
\end{quote}

\begin{quote}
\emph{sainike talvārthī ḍhāl vī̃dhī.}
\end{quote}

\begin{quote}
The soldier pierced the shield with his sword.
\end{quote}

Past, with \gu{તેણે} (he or she …): \gu{તેણે} \gu{ઢાલ} \gu{વીંધી}, \emph{teṇe ḍhāl vī̃dhī}.

\begin{cousinweb}
\textbf{\gu{વીંધવું}} (\emph{vī̃dhvũ}) comes from \emph{*vindhati}, a form CDIAL reconstructs, and CDIAL lists the Gujarati verb, \textquotedblleft{}to pierce\textquotedblright{} (CDIAL 11784). It agrees with what is pierced: \textbf{\gu{ઢાલ} \gu{વીંધી}}.
\end{cousinweb}

\begin{grammarlens}[title={What you've built}]
One more word for telling what happened.
\end{grammarlens}

\subsection*{Your turn}
\begin{itemize}
\raggedright
  \item \emph{Say it:} \emph{vī̃dhvũ}
  \item \emph{Say it:} \emph{vī̃dhvũ}, once more
  \item \emph{Say it:} \emph{sainike talvārthī ḍhāl vī̃dhī.} — The soldier pierced the shield with his sword.
  \item \emph{Recall:} say the Gujarati for a weapon, then the Gujarati for a cannon, then say \emph{vī̃dhvũ} again
\end{itemize}

\begin{tcolorbox}[breakable,colback=teal!4,colframe=teal!35!black,title={Before you move on}]
What does \gu{વીંધવું} mean? (\textquotedblleft{}To pierce\textquotedblright{}.) Say it once more.
\end{tcolorbox}
